\section{Introduction}
\label{sec:introduction}

The European Union Emissions Trading System (EU~ETS), operational since 2005, raises production costs for emission-intensive manufacturers through a cap on aggregate emissions and a market-determined permit price. Standard trade theory predicts that an asymmetric carbon cost shifts comparative advantage toward jurisdictions with weaker regulation and generates \emph{carbon leakage}: an increase in foreign emissions that partially offsets domestic abatement \citep{fowlie2016}. The incentive is strongest for homogeneous commodities traded at common world prices, where regulated producers cannot pass the carbon cost on to buyers and imports substitute for domestic output at the margin. Whether this mechanism operates through trade flows, as opposed to physical relocation, remains an open empirical question. Aggregate studies of EU manufacturing trade find no statistically significant leakage effect \citep{naegele2019}, yet aggregate estimates may mask leakage concentrated in specific bilateral corridors where regulatory asymmetry is large, transport costs are low, and production capacity in emission-intensive sectors already exists.

This paper asks one question: did increases in the effective EU~ETS carbon cost, net of free allocation, raise EU imports of emission-intensive goods from Morocco, Algeria, Egypt, and Tunisia relative to less emission-intensive goods over 2005 to 2021? It proposes and defends a design that would answer it. A positive coefficient would be consistent with the import-leakage channel operating in this corridor. It would not by itself establish physical relocation of production. North Africa is a natural candidate for corridor-specific leakage. The four partners hold Euro-Mediterranean association agreements that established free trade with the EU in industrial goods \citep{eumed}. They host cement, iron and steel, and nitrogen-fertiliser capacity oriented toward export markets, and Egypt ranked fifth among world urea exporters in 2020 \citep{usgs2021}. None of the four priced carbon at any point in the sample period \citep{worldbank2026}. These are precisely the sectors and product categories covered by the Carbon Border Adjustment Mechanism (CBAM; Regulation (EU) 2023/956), which began its transitional phase in October 2023 \citep{cbam2023}. The proposed sample window ends in 2021, before the CBAM proposal could plausibly affect trade patterns, so the design would quantify the pre-CBAM baseline that the mechanism is intended to correct.

The identification strategy rests on a triple-difference specification. Treatment intensity varies across products through their carbon-cost exposure, measured as the CO$_2$ emission intensity of the supplying sector. Treatment varies over time through the effective carbon price, defined as the annual European Union Allowance (EUA) price attenuated by the sector-year free-allocation share. Geographic variation distinguishes the North African corridor from a comparison partner group. Product-by-year fixed effects absorb global product-specific demand and supply shocks, such as Chinese steel overcapacity or commodity-price movements, that are correlated with emission intensity. Partner-by-year fixed effects absorb all time-varying partner-level confounds, including macroeconomic conditions and exchange-rate movements. The triple interaction is therefore identified from the differential response of high-exposure products sourced specifically from lower-regulation partners.

The paper contributes to the empirical literature on carbon leakage in three respects. First, it shifts the unit of analysis from EU-wide averages to a geographically targeted corridor where leakage incentives are strongest. This shift addresses the aggregation concern raised by \citet{verde2020}. Second, it incorporates free-allocation attenuation into the treatment definition, which the firm-level literature typically treats as a robustness dimension. The policy itself used free allocation to mitigate leakage risk \citep{martin2016}. Third, it proposes a pre-CBAM baseline estimate against which the mechanism's effectiveness could eventually be assessed.

The remainder of the paper is organised as follows. Section~\ref{sec:policy} sets out the institutional background insofar as it bears on the treatment definition. Section~\ref{sec:literature} reviews the empirical evidence by mechanism. Section~\ref{sec:design} presents the specification, its identifying assumptions, and the threat each answers. Section~\ref{sec:contribution} discusses what other researchers could take from the design, and Section~\ref{sec:conclusion} concludes.
